\section*{Computational interface and remaining proof obligations}
\addcontentsline{toc}{section}{Computational interface and remaining proof obligations}
\label{sec:computational-interface}

The computational layer encodes candidate data, guards, local transitions
and residual diagnostics. A coded instruction set specifies what a task
would mean; it does not establish that the required witness exists,
that a physical evolution realizes the instruction, or that a search
terminates with a closure certificate.

The complete local Zuse core, conditional universality obligations,
tool ledger, decoder and pre-OP status records are preserved in
Appendix~\ref{app:computational-protocols}. Historical declarations of
groundwork completion there concern the stated interface level, not
the unresolved OP~1/4 correction or OP~5 balance. Their use remains
subject to the hypotheses of the cited statements.

For a substantive completion, one still needs an explicit admissible
realization, a well-defined residual with its limiting/regularization
conventions, quantitative bounds and a witness satisfying the target
condition. Neither routing notation nor a successful document build
discharges these obligations.
